\rok{Децембар 2025.}{B}

Први практични тест из Алгоритама и структура података

\zadatak{1}{Quick Sort}
Написати \textit{quick} методу \textit{Quick Sort} алгоритма за сортирање целобројног низа дужине $n$.

\textbf{Додатне напомене за решавање задатка:}
\begin{itemize}
	\item Улаз је несортиран низ целих бројева.
	\item Излаз је сортирани низ целих бројева.
	\item У \textit{template}-у се налази класа \texttt{RandomNumberGen} коју треба искористити за генерисање низа случајних бројева.
	\item У оквиру класе \texttt{QuickSort} написати имплементацију за методу.
	\item Обезбедити код у \texttt{Main} методи за тестирање алгоритма.
\end{itemize}

\textbf{Потпис методе:}
\begin{Verbatim}[fontsize=\small]
private static void quickSort(int[] arr, int lower, int upper)
\end{Verbatim}

\textbf{Примери:}
Нису наведени у оригиналном PDF-у.

\zadatak{2}{Брисање непарних бројева из стека, производ преосталих}
Имплементирати методу која прима низ целих бројева. Метода треба да користи структуру података стека (класа \texttt{Stack}) за организацију и манипулацију елементима, и треба да изврши следеће кораке:
\begin{itemize}
	\item Уклањање непарних бројева: пролазом кроз стек, уклонити све бројеве који су непарни.
	\item Рачунање производа преосталих бројева: израчунати производ свих бројева који су остали у стеку након уклањања непарних бројева.
	\item Враћање резултата: као резултат, метода треба да врати израчунати производ преосталих бројева у стеку.
\end{itemize}

\textbf{Додатни захтеви за решавање задатка:}
\begin{itemize}
	\item Задатак решавати помоћу структуре стека (класа \texttt{Stack}) која је дата у оквиру шаблона задатка.
	\item Задатак решавати у оквиру методе \texttt{public int RemoveOddAndReturnProduct (int[] array)}.
	\item Након што се уклоне непарни бројеви из стека и израчуна производ преосталих бројева, стек треба да остане у истом стању као пре почетног процеса. То значи да преостали бројеви морају бити враћени у стек у оригиналном редоследу, а сви измењени елементи (они који су уклоњени) не смеју се вратити.
	\item Алгоритам је неопходно написати са линеарном временском комплексношћу $O(n)$.
\end{itemize}

\textbf{Потпис методе:}
\begin{Verbatim}[fontsize=\small]
public int RemoveOddAndReturnProduct (int[] array)
\end{Verbatim}

\textbf{Примери:}
\begin{itemize}
	\item \textbf{Улаз:} \texttt{[10, 15, 7, 30, 2, 8]} \\
		  \textbf{Излаз:} Производ преосталих бројева: \texttt{4800}; преостали бројеви у стеку: \texttt{10, 30, 2, 8}
	\item \textbf{Улаз:} \texttt{[5, 2, 8, 3, 4, 6]} \\
		  \textbf{Излаз:} Производ преосталих бројева: \texttt{384}; преостали бројеви у стеку: \texttt{2, 8, 4, 6}
\end{itemize}

\zadatak{3}{Josephus Problem}
Класичан \textit{Josephus problem} је математичка загонетка и алгоритамска вежба која потиче из античке легенде. Имате $n$ елемената (бројева) распоређених у круг (циркуларну листу). Елиминација се врши тако што се почне од одређеног места и броји се $M$ елемената у кругу, при чему се $M$-ти елемент елиминише из круга. Бројање се наставља од следећег елемента после елиминисаног. Процес се понавља док не остане само један елемент - „преживели“.

\textbf{Проблем:}
У класичном проблему, $M$ је константно (нпр. сваки трећи се елиминише). У нашем случају, $M$ се мења после сваке елиминације према формули:
$$
M_{novo} = (M_{staro} \times 2) \% N_{preziveli} + 1
$$

\textbf{Додатни захтеви за решавање задатка:}
\begin{itemize}
	\item Задатак решавати у методи \texttt{public int JosephusProblem(int startM)} која се налази у оквиру класе \texttt{CircularList}.
	\item Параметар \texttt{int startM} који се прослеђује методи \texttt{JosephusProblem} задаје корисник приликом првог позива ове методе.
	\item Сваки чвор листе треба да представља једног преживелог са својим ID бројем (од 1 до $n$).
	\item Обезбедити код у \texttt{Main} методи за тестирање.
\end{itemize}

\textbf{Потпис методе:}
\begin{Verbatim}[fontsize=\small]
public int JosephusProblem(int startM)
\end{Verbatim}

\textbf{Пример:}
\begin{itemize}
	\item \textbf{Улаз} (циркуларна листа): \texttt{1 -> 2 -> 3 -> 4 -> 5 -> 6}
	\item \textbf{startM:} \texttt{2}
	\item \textbf{Корак 1:} \texttt{1 -> 3 -> 4 -> 5 -> 6}; бројање стартује од: \texttt{3}; \texttt{novoM: 5}
	\item \textbf{Корак 2:} \texttt{3 -> 4 -> 5 -> 6}; бројање стартује од: \texttt{3}; \texttt{novoM: 3}
	\item \textbf{Корак 3:} \texttt{3 -> 4 -> 6}; бројање стартује од: \texttt{6}; \texttt{novoM: 1}
	\item \textbf{Корак 4:} \texttt{3 -> 4}; бројање стартује од: \texttt{3}; \texttt{novoM: 1}
	\item \textbf{Корак 5 (преживели):} \texttt{4}
\end{itemize}

\sablon{AISP2025_Test1_GrupaB}{2025/Decembar/GrupaB/AISP2025_Test1_GrupaB.linq}
